% ==========================================================
% SLIDE 48A: 2025 — From Idea to the Right Problem
% ==========================================================
\begin{frame}
    \begin{tikzpicture}[remember picture,overlay]
        \fill[PrimaryBlueLighter, opacity=0.08] ([xshift=-3cm,yshift=2cm]current page.south east) circle (4cm);
        \fill[PrimaryGoldLight, opacity=0.06] ([xshift=4cm,yshift=-2cm]current page.north west) circle (3.5cm);
        \node[anchor=north west, text=FundoQuestion, font=\bfseries\tamanhoQuestion\selectfont]
            at ([xshift=0.8cm,yshift=-0.8cm]current page.north west) {How did the idea become the right problem?};
        \node[anchor=north west, text width=.80\paperwidth, align=left, text=FundoKeyMessage, font=\tamanhoKeyMessage\selectfont]
            at ([xshift=0.8cm,yshift=-1.6cm]current page.north west) {
            We started by trying to build another Scan \& Go solution. By questioning the app, the customer journey, and the retailer's needs, we discovered that the real opportunity was solving the gap between physical retail and digital data. 
        };
    \end{tikzpicture}

    \vspace{0.2cm}
    \begin{center}
        {\Huge\bfseries\textcolor{PrimaryBlue}{From Idea to the Right Problem}}
        \vspace{0.05cm}
    \end{center}

    \vspace{-0.90cm}
    \begin{center}
        \begin{tikzpicture}[
            milestone/.style={
                draw=PrimaryBlue!30,
                rounded corners=4pt,
                align=left,
                fill=BlueCard,
                line width=0.8pt,
                font=\small,
                text width=14cm
            },
            friction/.style={
                draw=PrimaryBlue!30,
                rounded corners=3pt,
                minimum width=3.2cm,
                minimum height=0.7cm,
                align=left,
                fill=white,
                line width=0.5pt,
                font=\tiny
            },
            rproblem/.style={
                draw=PrimaryBlue!30,
                rounded corners=4pt,
                align=left,
                fill=BlueCard,
                line width=0.8pt,
                font=\small,
                text width=14cm
            }
        ]
            % ==================== PART 1: MILESTONES ====================
            \node[milestone, fill=BlueCard] (idea) at (-4.0, 6.22) {
                \textbf{Initial Idea (BB-Buy in a Blink)}: Turn an obsolete smartphone into a smart selling point, then take the solution to large physical retailers.
            };
            
            \node[milestone, fill=PrimaryGoldLight, draw=PrimaryGold] (questions) at (-4, 5.2) {
                \textbf{The Question:} Why build another Scan \& Go app when delivery platforms like iFood could do it?
            };
            
            \node[milestone, fill=GoldCard, draw=PrimaryGold] (insight) at (-4, 4.2) {
                \textbf{The Insight:} The core problem wasn't scanning. It was the gap between physical retail and digital data.
            };
             
            % ==================== PART 2: REAL PROBLEM TEXT ====================
            \node[rproblem, fill=GoldCard, draw=PrimaryGold] (insight) at (-4, 3.06) {

                \textbf{\textcolor{PrimaryRed}{The Real Problem:}} Price tags are static and manual, creating operational costs and errors. 
                They ultimately limit dynamic pricing. Our initial solution, however, required too many steps.

            };

% ==================== PART 3: FRICTION DIAGRAM ====================

            % Placing nodes manually to avoid matrix errors
            \node[friction, fill=BlueCard] (barcode) at (-9, 2.0) {Open a Barcode Reader};
            \node[friction, fill=BlueCard] (reader) at (-4, 2.0) {Scan the product code};
            \node[friction, fill=BlueCard] (copy) at (1, 2.0) {Copy the product code};

            \node[friction, fill=GoldCard] (whatsapp) at (-9, 1.0) {Paste and send the code};
            \node[friction, fill=GoldCard] (search) at (-4, 1.0) {Find the retailer's chat};
            \node[friction, fill=GoldCard] (send) at (1, 1.0) {Switch to WhatsApp};

            % Top-row flow
            \draw[->, thick, PrimaryBlue] (barcode.east) -- (reader.west);
            \draw[->, thick, PrimaryBlue] (reader.east) -- (copy.west);

            % Small vertical connection
            \draw[->, thick, PrimaryGold] (copy.south) -- (send.north);

            % Bottom-row flow
            \draw[->, thick, PrimaryGold] (send.west) -- (search.east);
            \draw[->, thick, PrimaryGold] (search.west) -- (whatsapp.east);

        \end{tikzpicture}
    \end{center}
\end{frame}